\begin{frame}[t]{Como SRAM e DRAM armazenam um bit}
  \begin{block}{SRAM: célula biestável (semelhante a um flip-flop)}
    Usa transistores realimentados para manter o estado 0 ou 1.
    A célula típica tem seis transistores e funciona como um
    \textit{latch}, sem precisar de clock ou de \textit{refresh}.
  \end{block}

  \begin{block}{DRAM: capacitor e transistor de acesso}
    Armazena o bit como carga elétrica em um capacitor.
    Como essa carga diminui com o tempo, precisa de
    \textit{refresh}: atualização periódica para preservar os dados.
  \end{block}

  \begin{alertblock}{As duas precisam de alimentação}
    Sem alimentação, ambas perdem os dados armazenados.
  \end{alertblock}
\end{frame}
